Das Boot fährt vor dem Wind, das Segel steht senkrecht zur Windrichtung. Für die Kraft auf das Segel zählt die Geschwindigkeit der Luft relativ zum Segel:
\begin{align}
 v_\text{rel} &= \vReF = \vWd-\vBt = \vRe
\end{align}
Da das Boot mit konstanter Geschwindigkeit fährt, ist die resultierende Kraft null: Der Wasserwiderstand ist gleich gross wie die Windkraft auf das Segel. Mit $c_w = \cwPO$ für eine quadratische Platte (FoTaBe S.~189) und $\rho = \rhLO$ (FoTaBe S.~187) ist
\begin{align}
 F_\text{Wasser} &= \FWaF \\
   &= \frac{1}{2}\cdot\cwP\cdot\rhL\cdot\ASe\cdot\left(\vRe\right)^2 \\
   &= \FWa \approx \result{\FWaP{0}{2}}
\end{align}
